\chapter{Sistemul real-time de procesare audio}

\section{Rolul sistemului}

Sistemul real-time de procesare audio reprezinta componenta interactiva a
proiectului de \textit{speech enhancement}. Daca pipeline-ul offline permite
antrenarea, evaluarea si compararea metodelor pe fisiere audio preinregistrate,
modul real-time are rolul de a demonstra comportamentul intregului sistem in
conditii de utilizare apropiate de o aplicatie practica: semnalul este captat
continuu, este procesat pe blocuri scurte, apoi este redat imediat catre
dispozitivul de iesire.

Obiectivul principal al acestei componente este reducerea zgomotului dintr-un
semnal de vorbire cu o intarziere suficient de mica pentru monitorizare live.
Aplicatia primeste un flux audio de la microfon sau dintr-o sursa sintetica
\texttt{speech+noise}, aplica una dintre metodele de imbunatatire disponibile
si livreaza catre difuzoare sau casti semnalul procesat. In paralel, interfata
grafica afiseaza formele de unda, spectrogramele, nivelurile RMS, scorul VAD
si statisticile de procesare.

Din punct de vedere functional, sistemul real-time trebuie sa indeplineasca
urmatoarele cerinte:

\begin{itemize}
    \item sa citeasca semnal audio in mod continuu, fara blocarea firului audio;
    \item sa proceseze fiecare bloc folosind o metoda selectabila: MaskNet,
    Spectral Subtraction, Wiener Filter sau bypass;
    \item sa redea semnalul imbunatatit cu latenta mica si cu amplitudine
    controlata;
    \item sa permita schimbarea metodei de procesare fara repornirea aplicatiei;
    \item sa ofere o interfata vizuala pentru diagnostic si prezentare;
    \item sa salveze optional inregistrari raw/enhanced si snapshot-uri ale
    interfetei.
\end{itemize}

Aceasta componenta este importanta deoarece face legatura dintre rezultatele
teoretice si utilizarea efectiva a modelului. Un model poate obtine valori bune
pentru PESQ, STOI sau SNR in evaluarea offline, dar in regim real-time mai
apar constrangeri suplimentare: timpul de inferenta, stabilitatea stream-ului
audio, alegerea dispozitivelor, tratarea blocurilor lipsa si controlul
feedback-ului acustic.

\section{Structura modulara}

Implementarea este impartita in mai multe module, fiecare avand o responsabilitate
clara. Separarea este utila deoarece partea de streaming audio, partea de
procesare, modelul neuronal si interfata grafica au cerinte diferite de timp,
sincronizare si testare.

\begin{table}[htbp]
\centering
\caption{Modulele principale ale sistemului real-time.}
\label{tab:module_realtime}
\vspace{0.5em}
\begin{tabular}{p{0.35\textwidth}p{0.57\textwidth}}
\toprule
\textbf{Modul} & \textbf{Responsabilitate} \\
\midrule
\path{src/tools/realtime_demo.py} &
Punctul de intrare al aplicatiei. Citeste argumentele CLI, alege checkpoint-ul,
configureaza dispozitivul PyTorch, dimensiunea blocurilor audio, dispozitivele
de intrare/iesire si porneste motorul real-time. \\
\addlinespace
\path{src/tools/realtime/engine.py} &
Contine clasele \texttt{StreamingEnhancer} si \texttt{RealtimeDenoiseApp}.
Gestioneaza bufferul de context, procesarea pe blocuri, cozile audio, firul
worker si comunicarea cu interfata. \\
\addlinespace
\path{src/tools/realtime/ui.py} &
Contine clasa \texttt{LiveComparisonUI}. Construieste interfata Tkinter,
controalele, graficele Matplotlib, spectrogramele, nivelurile RMS si curba VAD. \\
\addlinespace
\path{src/tools/realtime/common.py} &
Include functii comune: conversia milisecunde-esantioane, normalizarea
blocurilor, calculul RMS, alegerea dispozitivelor, recording-ul si sursa
sintetica \texttt{speech+noise}. \\
\addlinespace
\path{src/models/inference.py} &
Incarca modelul MaskNet din checkpoint si aplica inferenta pe un semnal audio,
folosind STFT, estimarea mastii si reconstructia prin ISTFT. \\
\addlinespace
\path{src/dsp/} &
Contine metodele clasice de procesare: Spectral Subtraction, Wiener Filter,
VAD Energy-ZCR si utilitarele STFT. \\
\bottomrule
\end{tabular}
\end{table}

Aceasta organizare separa clar trei planuri ale aplicatiei: planul audio
real-time, planul de procesare si planul vizual. Planul audio trebuie sa fie
cat mai rapid si previzibil; planul de procesare poate fi mai costisitor, dar
este izolat intr-un fir separat; planul vizual poate fi actualizat cu o rata
mai mica decat rata callback-urilor audio, deoarece utilizatorul nu are nevoie
de redesenare la fiecare bloc.

\section{Parametrii audio de baza}

Sistemul foloseste aceeasi configuratie audio ca restul proiectului, pentru a
pastra compatibilitatea dintre antrenare, evaluare, demo offline si demo
real-time. Frecventa de esantionare este:

\[
f_s = 16000 \ \text{Hz}.
\]

Parametrii STFT sunt:

\begin{itemize}
    \item lungimea ferestrei: \(L = 400\) esantioane, adica \(25\) ms;
    \item hop-ul: \(H = 160\) esantioane, adica \(10\) ms;
    \item dimensiunea FFT: \(N_{\text{FFT}} = 512\);
    \item numarul de benzi de frecventa: \(K = 257\).
\end{itemize}

Pentru streaming, dimensiunea implicita a blocului audio este de \(32\) ms:

\[
B = \left\lfloor f_s \cdot \frac{32}{1000} \right\rfloor = 512
\quad \text{esantioane}.
\]

In mod implicit, procesarea foloseste si o fereastra de context de \(640\) ms:

\[
C = \left\lfloor f_s \cdot \frac{640}{1000} \right\rfloor = 10240
\quad \text{esantioane}.
\]

Contextul nu reprezinta in intregime latenta audibila. El este o memorie
istorica folosita pentru ca algoritmii bazati pe STFT si modelul MaskNet sa
lucreze pe o secventa suficient de lunga. Latenta efectiva perceputa este
influentata mai ales de dimensiunea blocului, de intarzierea de taiere a cozii,
de timpul de procesare si de buffering-ul dispozitivului audio.

\section{Modelul de semnal}

Semnalul observat la intrare este modelat ca suma dintre vorbirea curata si
zgomot:

\[
x[n] = s[n] + v[n],
\]

unde \(s[n]\) este semnalul curat de vorbire, \(v[n]\) este zgomotul, iar
\(x[n]\) este semnalul captat de microfon sau obtinut prin mixarea unui fisier
de vorbire cu unul de zgomot.

Scopul sistemului de imbunatatire este estimarea unui semnal:

\[
\hat{s}[n] \approx s[n],
\]

prin reducerea componentelor asociate zgomotului si pastrarea cat mai buna a
continutului vocal. In domeniul timp-frecventa, semnalul este reprezentat prin
STFT:

\[
X[k,m] =
\sum_{n=0}^{L-1} x[n + mH] \, w[n] \, e^{-j2\pi kn/N_{\text{FFT}}},
\]

unde \(k\) este indicele frecventei, \(m\) este indicele cadrului, \(w[n]\)
este fereastra Hann, \(L\) este lungimea cadrului, iar \(H\) este hop-ul.

Cele mai multe metode din proiect opereaza asupra magnitudinii:

\[
|X[k,m]|,
\]

pastrand faza semnalului zgomotos pentru reconstructie. In varianta complexa a
modelului MaskNet, sistemul poate lucra si cu partea reala si imaginara a STFT,
estimand o masca complexa.

\section{Fluxul general de executie}

La pornire, aplicatia real-time urmeaza o secventa clara de initializare:

\begin{enumerate}
    \item sunt citite argumentele din linia de comanda;
    \item este selectata metoda de procesare: \texttt{masknet},
    \texttt{spectral\_subtraction}, \texttt{wiener} sau \texttt{bypass};
    \item daca metoda este \texttt{masknet}, este ales si incarcat checkpoint-ul;
    \item este selectat dispozitivul PyTorch, de obicei \texttt{cuda} daca exista
    GPU disponibil, altfel \texttt{cpu};
    \item sunt convertiti parametrii temporali din milisecunde in esantioane;
    \item este creat obiectul \texttt{StreamingEnhancer};
    \item sunt alese dispozitivele audio de intrare si iesire;
    \item optional, este creata interfata \texttt{LiveComparisonUI};
    \item optional, este creata sesiunea de recording;
    \item este pornita aplicatia \texttt{RealtimeDenoiseApp}.
\end{enumerate}

Comanda uzuala de pornire a demo-ului cu interfata grafica este:

\begin{lstlisting}[language=bash]
python -m src.tools.project_cli realtime-demo -- --method masknet --visualize
\end{lstlisting}

Pentru verificarea dispozitivelor audio, se poate folosi:

\begin{lstlisting}[language=bash]
python -m src.tools.project_cli realtime-demo -- --list-devices
\end{lstlisting}

Iar pentru o demonstratie cu salvarea semnalelor raw si enhanced:

\begin{lstlisting}[language=bash]
python -m src.tools.project_cli realtime-demo -- --method masknet --visualize --record
\end{lstlisting}

\section{Captura si redarea audio}

Fluxul audio este realizat cu biblioteca \texttt{sounddevice}. Implementarea
foloseste doua stream-uri separate: un \texttt{InputStream} pentru citirea
semnalului de la microfon si un \texttt{OutputStream} pentru redarea semnalului
procesat. Ambele ruleaza cu acelasi sample rate, aceeasi dimensiune de bloc si
un singur canal audio.

In callback-ul de intrare, sistemul copiaza blocul primit de la dispozitiv:

\[
x_i = [x_i[0], x_i[1], \ldots, x_i[B-1]],
\]

unde \(i\) este indicele blocului curent, iar \(B\) este dimensiunea blocului
in esantioane. Daca sursa \texttt{speech+noise} este activa si exista fisiere
incarcate, blocul de microfon este inlocuit cu un bloc extras din mixul
sintetic.

Callback-ul de intrare nu ruleaza algoritmul de denoising. El doar pune blocul
in coada de intrare:

\[
x_i \rightarrow \texttt{input\_queue}.
\]

Callback-ul de iesire incearca sa extraga cel mai recent bloc procesat:

\[
\hat{s}_i \leftarrow \texttt{output\_queue}.
\]

Daca nu exista un bloc disponibil, sistemul trimite la iesire un bloc de
liniste. Aceasta alegere este preferabila blocarii callback-ului audio, deoarece
un callback blocat produce intreruperi audibile si instabilitate la nivelul
driverului audio.

\section{Modelul concurent de procesare}

Procesarea este separata de callback-urile audio printr-un model
\textit{producer-consumer}. Callback-ul de intrare produce blocuri raw, firul
worker le consuma si produce blocuri enhanced, iar callback-ul de iesire consuma
blocurile enhanced.

Sistemul foloseste doua cozi cu dimensiune limitata:

\begin{itemize}
    \item \texttt{input\_queue}: coada blocurilor raw;
    \item \texttt{output\_queue}: coada blocurilor procesate.
\end{itemize}

Firul worker ruleaza continuu metoda de procesare:

\begin{lstlisting}[language=Python]
block = self.input_queue.get(timeout=0.1)
enhanced = self.enhancer.process_block(block)
self.output_queue.put_nowait(enhanced)
\end{lstlisting}

Acest design are doua avantaje. In primul rand, callback-urile audio raman
scurte si predictibile. In al doilea rand, metodele costisitoare precum MaskNet
pot rula fara a bloca direct firul audio. Daca procesarea este temporar mai
lenta decat ritmul audio, coada de iesire poate ramane goala, caz in care se
redau blocuri de liniste. Daca intrarea se acumuleaza prea mult, sistemul
elimina blocuri vechi pentru a evita cresterea necontrolata a latentei.

Oprirea este controlata prin \texttt{threading.Event}. Atat bucla principala,
cat si firul worker verifica periodic evenimentul de oprire:

\begin{lstlisting}[language=Python]
while not self.stop_event.is_set():
    ...
\end{lstlisting}

La inchiderea ferestrei, la apasarea butonului \texttt{Stop Demo} sau la
\texttt{Ctrl+C}, evenimentul este setat, stream-urile sunt inchise, firul worker
este asteptat prin \texttt{join}, iar fisierele de recording sunt salvate daca
exista date.

\section{Procesarea cu buffer de context}

Clasa \texttt{StreamingEnhancer} este responsabila pentru aplicarea metodei de
denoising pe blocuri audio. Deoarece algoritmii bazati pe STFT si modelul
MaskNet au nevoie de context temporal, fiecare bloc nu este procesat complet
izolat. In schimb, sistemul pastreaza un buffer glisant:

\[
u_i =
[x_{i-r}, x_{i-r+1}, \ldots, x_i],
\]

unde \(u_i\) este fereastra curenta, iar \(r\) este numarul de blocuri istorice
pastrate in context. In implementare, dimensiunea totala a bufferului este cel
putin:

\[
N_{\text{buffer}} = C + T + B,
\]

unde \(C\) este contextul, \(T\) este zona de coada taiata pentru reducerea
artefactelor de margine, iar \(B\) este blocul curent.

La fiecare pas, bufferul este deplasat spre stanga, iar blocul nou este copiat
la final:

\[
u_i[n] =
\begin{cases}
u_{i-1}[n+B], & 0 \leq n < N_{\text{buffer}} - B, \\
x_i[n - (N_{\text{buffer}} - B)], & N_{\text{buffer}} - B \leq n < N_{\text{buffer}}.
\end{cases}
\]

Metoda de denoising este aplicata pe intreaga fereastra \(u_i\), rezultand
\(\hat{u}_i\). Din aceasta fereastra procesata se extrage doar blocul care
trebuie redat:

\[
\hat{s}_i =
\hat{u}_i[N_{\text{buffer}} - T - B : N_{\text{buffer}} - T].
\]

Taierea cozii cu \(T\) esantioane reduce efectele de margine aparute la finalul
ferestrei procesate. In configuratia implicita, \(T\) este egal cu lungimea
ferestrei STFT, adica \(25\) ms. Astfel, sistemul prefera sa redea un segment
aflat putin in urma capatului curent al bufferului, deoarece acesta are mai mult
context si este mai stabil.

\section{Normalizare si controlul amplitudinii}

In procesarea real-time, amplitudinea trebuie controlata cu atentie pentru a
evita clipping-ul si semnalele instabile. Fiecare bloc audio este convertit in
\texttt{float32}, valorile invalide sunt inlocuite cu zero, iar semnalul este
limitat in intervalul \([-1, 1]\).

Functia de normalizare foloseste un prag de varf:

\[
p = \max_n |x[n]|.
\]

Daca \(p > 0.98\), blocul este scalat:

\[
x_{\text{norm}}[n] = x[n] \cdot \frac{0.98}{p}.
\]

Altfel, blocul este pastrat. La final, se aplica o limitare:

\[
x_{\text{out}}[n] = \min(1, \max(-1, x_{\text{norm}}[n])).
\]

Sistemul include si un control \texttt{dry/wet}, care permite amestecarea
semnalului procesat cu semnalul original:

\[
y[n] = \lambda \hat{s}[n] + (1-\lambda)x[n],
\]

unde \(\lambda \in [0,1]\). Pentru \(\lambda = 1\), utilizatorul aude doar
semnalul imbunatatit. Pentru \(\lambda = 0\), sistemul functioneaza ca monitor
direct al intrarii.

Exista si un control de \texttt{output gain}, exprimat in dB. Conversia in
factor liniar este:

\[
g = 10^{G_{\text{dB}}/20}.
\]

Semnalul redat devine:

\[
y_g[n] = g \cdot y[n],
\]

dupa care se aplica din nou normalizarea pentru siguranta.

\section{Metodele de procesare disponibile}

\subsection{Bypass}

Modul \texttt{bypass} nu aplica nicio filtrare. Semnalul de iesire este identic
cu intrarea, cu exceptia normalizarii de siguranta:

\[
\hat{s}[n] = x[n].
\]

Acest mod este util in demonstratii deoarece ofera o referinta auditiva si
vizuala. Utilizatorul poate asculta semnalul brut, apoi poate comuta catre
Spectral Subtraction, Wiener Filter sau MaskNet pentru a observa diferenta.

\subsection{Spectral Subtraction}

Metoda Spectral Subtraction presupune estimarea unei magnitudini medii a
zgomotului si scaderea acesteia din magnitudinea semnalului zgomotos. In
implementarea proiectului, zgomotul este estimat din cadrele cu cea mai mica
energie, fara a necesita etichete externe VAD.

Pentru spectrograma zgomotoasa \(Y[k,m]\), energia cadrului este:

\[
E[m] = \frac{1}{K}\sum_{k=0}^{K-1}|Y[k,m]|^2.
\]

Se aleg primele \(M\) cadre cu energie minima, iar profilul de zgomot este:

\[
\hat{N}[k] =
\frac{1}{M}\sum_{m \in \mathcal{I}_{\text{low}}} |Y[k,m]|.
\]

Magnitudinea imbunatatita este:

\[
|\hat{S}[k,m]| =
\max\left(|Y[k,m]| - \alpha \hat{N}[k], \beta \hat{N}[k]\right),
\]

unde \(\alpha\) este factorul de scadere, iar \(\beta\) este podeaua spectrala.
In configuratia implicita, \(\alpha = 1.0\), iar \(\beta = 0.02\).

Spectrograma complexa este reconstruita folosind faza zgomotoasa:

\[
\hat{S}[k,m] = |\hat{S}[k,m]| e^{j\angle Y[k,m]}.
\]

Avantajul metodei este costul redus de calcul. Limitarea principala este
tendinta de a produce artefacte muzicale atunci cand estimarea zgomotului este
imprecisa sau zgomotul variaza rapid.

\subsection{Wiener Filter}

Filtrul Wiener foloseste o estimare a densitatii spectrale de putere a zgomotului
si calculeaza un castig dependent de raportul semnal-zgomot a posteriori.
Profilul de zgomot este estimat tot din cadrele cu energie mica:

\[
\hat{\sigma}_v^2[k] =
\frac{1}{M}\sum_{m \in \mathcal{I}_{\text{low}}} |Y[k,m]|^2.
\]

SNR-ul a posteriori este:

\[
\text{SNR}_{\text{post}}[k,m] =
\frac{|Y[k,m]|^2}{\hat{\sigma}_v^2[k] + \epsilon}.
\]

Castigul Wiener este:

\[
G[k,m] =
\frac{\text{SNR}_{\text{post}}[k,m]}
{\text{SNR}_{\text{post}}[k,m] + 1}.
\]

Magnitudinea procesata devine:

\[
|\hat{S}[k,m]| = G[k,m]|Y[k,m]|.
\]

Ca si Spectral Subtraction, metoda pastreaza faza semnalului zgomotos pentru
ISTFT. In practica, Wiener Filter este stabil si rapid, dar poate lasa mai mult
zgomot rezidual decat un model neuronal antrenat pe exemple variate.

\subsection{MaskNet}

Metoda \texttt{masknet} foloseste reteaua neuronala antrenata in proiect.
Modelul primeste spectrograma semnalului zgomotos si estimeaza o masca spectrala.
Pentru varianta de magnitudine, intrarea este:

\[
|Y| \in \mathrm{R}^{K \times M}.
\]

Modelul produce:

\[
\hat{M}[k,m] = f_{\theta}(|Y|)[k,m],
\]

unde \(f_{\theta}\) reprezinta reteaua MaskNet, iar \(\theta\) sunt parametrii
invatati in timpul antrenarii. Magnitudinea imbunatatita este:

\[
|\hat{S}[k,m]| = \hat{M}[k,m] \cdot |Y[k,m]|.
\]

Reconstructia se face cu faza semnalului zgomotos:

\[
\hat{S}[k,m] =
|\hat{S}[k,m]|e^{j\angle Y[k,m]}.
\]

Arhitectura MaskNet este de tip U-Net encoder-decoder, cu skip connections,
blocuri convolutionale si, in variantele mai avansate, LSTM bidirectional si
attention pe canale~\cite{unet,lstm,se}. Encoderul reduce rezolutia
timp-frecventa si extrage reprezentari compacte, bottleneck-ul modeleaza
dependente temporale, iar decoderul reconstruieste o masca la aceeasi rezolutie
cu intrarea.

In varianta complexa, modelul primeste doua canale:

\[
Y_{\text{in}} =
\begin{bmatrix}
\Re(Y) \\
\Im(Y)
\end{bmatrix},
\]

si estimeaza o masca complexa:

\[
\hat{M}_{\mathrm{C}} =
\hat{M}_r + j\hat{M}_i.
\]

Spectrograma imbunatatita este:

\[
\hat{S}[k,m] =
\hat{M}_{\mathrm{C}}[k,m] \cdot Y[k,m].
\]

Aceasta formulare permite corectii partiale si asupra fazei, nu doar asupra
magnitudinii~\cite{complex_mask}. In modul real-time, checkpoint-ul incarcat
determina automat daca modelul este de magnitudine sau complex.

\section{Voice Activity Detection in sistemul real-time}

Detectorul VAD folosit in proiect este bazat pe energie, rata trecerilor prin
zero si caracteristici spectrale. In interfata real-time, VAD-ul este folosit
in principal pentru monitorizare si diagnostic vizual: utilizatorul vede o
curba de incredere in activitatea vocala, calculata pe bufferul de intrare.

Pentru un cadru \(m\), energia este:

\[
E[m] = \frac{1}{L}\sum_{n=0}^{L-1}x_m^2[n].
\]

Rata trecerilor prin zero este:

\[
\text{ZCR}[m] =
\frac{1}{L-1}
\sum_{n=0}^{L-2}
\mathbf{1}\{\operatorname{sign}(x_m[n]) \neq
\operatorname{sign}(x_m[n+1])\}.
\]

Detectorul modernizat mai calculeaza raportul de energie in banda de vorbire
si flatness-ul spectral. Scorurile normalizate sunt combinate intr-un scor
soft:

\[
q[m] =
w_E q_E[m] +
w_B q_B[m] +
w_Z q_Z[m] +
w_F q_F[m],
\]

unde \(q_E\) este scorul de energie, \(q_B\) scorul de banda de vorbire,
\(q_Z\) scorul ZCR, iar \(q_F\) scorul de flatness. Ponderile sunt normalizate,
iar scorul final este limitat in intervalul \([0,1]\).

Decizia hard foloseste praguri cu histerezis:

\[
\text{VAD}[m] =
\begin{cases}
1, & q[m] \geq \theta_{\text{start}}, \\
0, & q[m] < \theta_{\text{end}},
\end{cases}
\quad
\theta_{\text{start}} > \theta_{\text{end}}.
\]

Sistemul foloseste si netezire temporala, eliminarea segmentelor prea scurte
si o forma de \textit{hangover}, astfel incat decizia sa nu oscileze rapid la
limita dintre vorbire si tacere.

In versiunea curenta a demo-ului, VAD-ul din interfata nu inlocuieste metoda de
enhancement si nu opreste direct procesarea audio. Rolul lui principal este de
a arata utilizatorului unde sistemul detecteaza activitate vocala si de a
sprijini interpretarea spectrogramelor.

\section{Sursa sintetica speech+noise}

Pe langa intrarea de microfon, aplicatia poate folosi o sursa sintetica numita
\texttt{speech+noise}. Aceasta permite incarcarea unui fisier de vorbire si a
unui fisier de zgomot, apoi mixarea lor la un SNR controlabil. Modul este util
pentru prezentari, deoarece reduce dependenta de zgomotul real din camera.

La incarcarea unui fisier audio, sistemul:

\begin{itemize}
    \item citeste fisierul cu \texttt{soundfile};
    \item converteste semnalele stereo la mono prin medierea canalelor;
    \item resampleaza semnalul la \(16\) kHz daca este necesar;
    \item normalizeaza amplitudinea pentru siguranta;
    \item reconstruieste mixul daca exista si fisier de zgomot incarcat.
\end{itemize}

Pentru mixarea la un SNR dorit, se calculeaza puterile medii:

\[
P_s = \frac{1}{N}\sum_{n=0}^{N-1}s^2[n],
\qquad
P_v = \frac{1}{N}\sum_{n=0}^{N-1}v^2[n].
\]

Daca SNR-ul dorit este \(R_{\text{dB}}\), factorul de scalare al zgomotului este:

\[
a =
\sqrt{
10^{-R_{\text{dB}}/10}
\cdot
\frac{P_s}{P_v + \epsilon}
}.
\]

Mixul rezultat este:

\[
x[n] = s[n] + a v[n].
\]

Sursa sintetica include si distorsiuni optionale, reutilizand pipeline-ul de
augmentare:

\begin{itemize}
    \item modificare de gain;
    \item time stretch;
    \item pitch shift;
    \item reverberatie;
    \item egalizare;
    \item compresie dinamica;
    \item clipping.
\end{itemize}

Interfata include preset-uri de scenariu. \texttt{neutral} foloseste un caz
moderat, \texttt{office} adauga EQ si compresie, \texttt{street} foloseste un
SNR mai mic si distorsiuni mai agresive, iar \texttt{hard} combina mai multe
distorsiuni pentru a testa limitele sistemului.

\section{Interfata grafica}

Interfata este construita cu Tkinter, \texttt{ttk} si Matplotlib. Fereastra este
impartita in doua zone principale. In stanga sunt afisate datele live, iar in
dreapta se afla controalele.

Zona de monitorizare include:

\begin{itemize}
    \item forma de unda a semnalului raw;
    \item forma de unda a semnalului enhanced;
    \item spectrograma semnalului raw;
    \item spectrograma semnalului enhanced;
    \item nivelurile RMS pentru intrare si iesire;
    \item indicatorul de semnal activ;
    \item scorul VAD soft;
    \item sumarul sesiunii.
\end{itemize}

Pentru afisarea formelor de unda, UI-ul pastreaza doua buffere circulare:
\texttt{raw\_buffer} si \texttt{output\_buffer}. La fiecare bloc nou, datele
vechi sunt deplasate, iar datele noi sunt inserate la final:

\[
b_i[n] =
\begin{cases}
b_{i-1}[n+B], & 0 \leq n < D-B, \\
x_i[n-(D-B)], & D-B \leq n < D,
\end{cases}
\]

unde \(D\) este numarul de esantioane afisat in istoric.

Spectrogramele sunt calculate local pentru afisare. Pentru fiecare frame se
aplica o fereastra Hann si FFT real, iar magnitudinea este convertita in dB:

\[
S_{\text{dB}}[k,m] =
20 \log_{10}\left(\max(|X[k,m]|, 10^{-4})\right).
\]

Valorile sunt limitate intre \(-80\) dB si \(0\) dB pentru stabilitate vizuala.
Interfata nu redeseneaza figura la fiecare callback audio, ci limiteaza rata de
redesenare la aproximativ 12 cadre pe secunda. Aceasta decizie reduce consumul
CPU si pastreaza UI-ul fluid.

\section{Controalele utilizatorului}

Controalele sunt grupate in tab-uri:

\begin{itemize}
    \item \textbf{Audio}: selectarea sursei \texttt{microphone} sau
    \texttt{speech+noise}, alegerea dispozitivelor de intrare si iesire;
    \item \textbf{Processing}: alegerea metodei de denoising, selectarea modului
    VAD, reglarea \texttt{dry/wet}, reglarea \texttt{output gain}, snapshot,
    recording si oprirea demo-ului;
    \item \textbf{Files}: incarcarea fisierelor de vorbire si zgomot, plus
    controlul play/pause pentru sursa sintetica;
    \item \textbf{Augment}: alegerea scenariului, reglarea SNR-ului si activarea
    distorsiunilor.
\end{itemize}

Comenzile din UI nu modifica direct starea motorului audio in momentul
evenimentului grafic. Ele seteaza variabile \texttt{pending}, iar bucla principala
le consuma intr-un punct controlat. Acest model se poate descrie astfel:

\[
\text{UI event} \rightarrow \text{pending request}
\rightarrow \text{engine consume} \rightarrow \text{state update}.
\]

Prin aceasta decizie, schimbarea metodei, a dispozitivelor, a sursei, a SNR-ului
sau a distorsiunilor nu se produce in interiorul callback-ului audio. Riscul de
conditii de cursa este redus, iar callback-ul ramane simplu.

\section{Recording si snapshot-uri}

Sistemul poate salva semnalele raw si enhanced in fisiere WAV. Cand recording-ul
este activ, fiecare pereche de blocuri este copiata in memorie:

\[
\{x_i\}_{i=1}^{N_b},
\qquad
\{\hat{s}_i\}_{i=1}^{N_b}.
\]

La finalul sesiunii, blocurile sunt concatenate:

\[
x_{\text{rec}} =
[x_1, x_2, \ldots, x_{N_b}],
\qquad
\hat{s}_{\text{rec}} =
[\hat{s}_1, \hat{s}_2, \ldots, \hat{s}_{N_b}],
\]

si sunt salvate ca:

\begin{itemize}
    \item \path{raw_microphone.wav};
    \item \path{enhanced_output.wav}.
\end{itemize}

Interfata poate salva si snapshot-uri PNG ale figurii principale, utile in
documentatie sau prezentare. Snapshot-ul surprinde forma de unda raw, forma de
unda enhanced si cele doua spectrograme in acelasi moment al sesiunii.

\section{Alegerea dispozitivelor audio}

Aplicatia poate lista dispozitivele audio disponibile si poate primi explicit
un ID sau un nume de dispozitiv. Daca utilizatorul nu alege manual dispozitivele,
sistemul incearca sa selecteze automat o pereche potrivita de intrare si iesire.

Selectia automata ia in calcul:

\begin{itemize}
    \item numarul de canale de intrare sau iesire;
    \item potrivirea sample rate-ului implicit cu \(16\) kHz;
    \item dispozitivele implicite ale sistemului;
    \item cuvinte-cheie precum \texttt{microphone}, \texttt{speaker},
    \texttt{headset}, \texttt{usb}, \texttt{bluetooth};
    \item penalizari pentru dispozitive virtuale sau nepotrivite.
\end{itemize}

Cand utilizatorul schimba dispozitivele din UI, stream-ul audio este repornit.
Daca noua pereche de dispozitive produce o eroare, aplicatia revine la ultima
pereche functionala. Aceasta masura este importanta deoarece driverele audio pot
refuza anumite combinatii de input si output, mai ales cand folosesc host API-uri
diferite.

\section{Robustete in regim real-time}

Aplicatia include mai multe mecanisme de robustete:

\begin{itemize}
    \item \textbf{Warmup}: inainte de pornirea stream-ului, enhancer-ul proceseaza
    o fereastra de liniste pentru a reduce latenta primului bloc;
    \item \textbf{Output underflow}: daca nu exista bloc procesat disponibil,
    se reda temporar liniste;
    \item \textbf{Input overflow}: daca intrarea se acumuleaza prea mult, se
    elimina blocuri vechi pentru a nu creste latenta;
    \item \textbf{Silent input warning}: daca microfonul ramane aproape silentios
    mai multe callback-uri, sistemul avertizeaza utilizatorul;
    \item \textbf{Device fallback}: daca o pereche noua de dispozitive nu poate
    fi deschisa, se revine la ultima pereche valida;
    \item \textbf{Controlled shutdown}: oprirea inchide stream-ul, thread-ul
    worker, UI-ul si salveaza recording-ul daca exista.
\end{itemize}

Aceste mecanisme sunt necesare deoarece demo-ul real-time este expus la factori
care nu apar in evaluarea offline: lipsa semnalului de microfon, incompatibilitati
de driver, procesare mai lenta decat ritmul audio, feedback acustic sau schimbari
rapide ale comenzilor UI.

\section{Analiza latentei}

Latenta totala a sistemului poate fi aproximata prin suma mai multor componente:

\[
T_{\text{total}} \approx
T_{\text{device}} +
T_{\text{block}} +
T_{\text{tail}} +
T_{\text{proc}} +
T_{\text{queue}},
\]

unde:

\begin{itemize}
    \item \(T_{\text{device}}\) este latenta driverului audio;
    \item \(T_{\text{block}}\) este durata unui bloc audio;
    \item \(T_{\text{tail}}\) este intarzierea introdusa prin taierea cozii;
    \item \(T_{\text{proc}}\) este timpul de procesare al metodei alese;
    \item \(T_{\text{queue}}\) este intarzierea introdusa de cozile interne.
\end{itemize}

Pentru valorile implicite:

\[
T_{\text{block}} = 32 \ \text{ms},
\qquad
T_{\text{tail}} = 25 \ \text{ms}.
\]

Contextul de \(640\) ms nu trebuie interpretat ca latenta audibila directa, ci
ca memorie de analiza. Totusi, el creste numarul de esantioane procesate la
fiecare apel si poate mari timpul de inferenta. De aceea, pe CPU, metoda MaskNet
poate fi mai dificil de rulat real-time decat metodele clasice. In astfel de
cazuri, aplicatia permite comutarea live la \texttt{spectral\_subtraction} sau
\texttt{wiener}.

Conditia practica pentru rulare stabila este:

\[
T_{\text{proc}} < T_{\text{block}}
\]

in medie, sau cel putin suficient de aproape de aceasta limita incat cozile sa
nu acumuleze intarzieri. Daca timpul de procesare depaseste constant durata
blocului, apar underflow-uri la iesire, drop-uri la intrare sau cresterea
latentei.

\section{Integrarea cu modelul antrenat}

Pentru metoda MaskNet, aplicatia incarca automat cel mai bun checkpoint local
daca utilizatorul nu furnizeaza explicit unul. Cautarea favorizeaza checkpoint-uri
recomandate, de exemplu experimentele LibriSpeech cu VAD soft sau variantele de
benchmark rapid, apoi cauta alte fisiere \path{masknet_best.pth}.

Checkpoint-ul contine atat parametrii modelului, cat si configuratia arhitecturii.
La incarcare, se reconstruiesc optiunile relevante:

\begin{itemize}
    \item numarul de canale de intrare;
    \item numarul de canale de iesire;
    \item numarul de canale de baza;
    \item folosirea LSTM;
    \item folosirea attention;
    \item folosirea blocurilor reziduale;
    \item tipul mastii: magnitudine sau complex;
    \item activarea de iesire: sigmoid, tanh sau liniara.
\end{itemize}

Modelul este trecut in modul \texttt{eval}, iar inferenta foloseste
\texttt{torch.inference\_mode}. Daca dispozitivul este CUDA, inferenta poate
folosi autocast pentru eficienta. Rezultatul este convertit inapoi in NumPy si
trimis catre motorul de streaming.

\section{Avantaje si limitari}

Principalul avantaj al sistemului real-time este ca pune toate componentele
proiectului intr-un flux demonstrabil: captare, procesare, redare, vizualizare,
control si salvare. Utilizatorul poate compara imediat metodele si poate vedea
efectul lor pe semnalul audio, nu doar in tabele de evaluare.

Avantajele principale sunt:

\begin{itemize}
    \item arhitectura modulara si usor de extins;
    \item suport pentru metode DSP rapide si pentru MaskNet;
    \item schimbare live a metodei de procesare;
    \item sursa sintetica \texttt{speech+noise} cu SNR controlabil;
    \item vizualizare simultana a formelor de unda si spectrogramelor;
    \item mecanisme de robustete pentru cozi, dispozitive si oprire;
    \item posibilitatea de a salva exemple audio si imagini pentru analiza.
\end{itemize}

Exista si limitari. MaskNet poate fi mai costisitor pe CPU, iar performanta
real-time depinde de hardware. Estimarea bazata pe ferestre poate produce
artefacte de margine daca dimensiunea contextului este prea mica. De asemenea,
metodele care pastreaza faza semnalului zgomotos sunt limitate atunci cand
distorsiunea de faza contribuie semnificativ la degradarea perceptuala. In
plus, utilizarea microfonului si a difuzoarelor laptopului poate produce
feedback acustic, de aceea castile sunt recomandate pentru demonstratii.

\section{Concluzie}

Sistemul real-time de procesare audio completeaza pipeline-ul de speech
enhancement printr-o componenta interactiva si demonstrabila. Arhitectura sa
se bazeaza pe separarea clara dintre callback-urile audio, firul de procesare
si interfata grafica. Blocurile audio sunt transportate prin cozi, procesarea
se face pe un buffer glisant cu context, iar metodele disponibile acopera atat
baseline-uri DSP, cat si modelul neuronal MaskNet.

Prin integrarea sursei \texttt{speech+noise}, a controalelor pentru SNR si
distorsiuni, a vizualizarii live si a functiilor de recording, aplicatia permite
testarea sistemului in conditii variate. Astfel, demo-ul real-time nu este doar
o interfata de prezentare, ci si un instrument de diagnostic pentru stabilitatea,
latenta si comportamentul perceptual al metodelor de reducere a zgomotului.
